% TOP_PROBLEM: {"id": 2800702, "title": "10 Lectures and 42 Open Problems — Boolean classification and annulus conjecture", "category_id": 15, "category": {"id": 15, "name": "computer_science", "display_name": "Computer Science", "description": "Computational complexity, algorithms, and theoretical CS.", "slug": "computer-science", "order_index": 15, "created_at": "2026-07-31T15:26:25.671Z"}, "status": "partially_solved"}
% FIRST_POSTED: null
% ENTRY_KIND: statement_only
% Imported from: batch08.tex; UnsolvedMath ID 2800702
% Compile twice: lualatex 2800702.tex
\documentclass[11pt]{article}
\usepackage{fontspec}
\setmainfont{Latin Modern Roman}
\usepackage{amsmath,amssymb,mathtools,unicode-math}
\setmathfont{Latin Modern Math}
\usepackage{xurl,hyperref}
\hypersetup{hidelinks,pdftitle={UnsolvedMath Problem 2800702}}
\newcommand{\sref}[2]{\hyperlink{source:#1:#2}{[S#2]}}
\newcommand{\cataloguescope}[1]{\paragraph{Mathematical question.}#1}
\begin{document}
% UnsolvedMath ID 2800702; source catalogue code AMR-027-0702
\setcounter{section}{2800701}
\section{10 Lectures and 42 Open Problems — Boolean classification and annulus conjecture}\label{2800702}
{\small\sffamily UnsolvedMath ID 2800702 · Computer Science}
\subsection{Definitions and mathematical statement}

\paragraph{Shared notation.}$\mathbb N=\{1,2,\ldots\}$, and $\mathbb Z,\mathbb Q,\mathbb R,\mathbb C$ denote the integers, rationals, reals and complex numbers. Any other conventions are specified locally.


Write \(\mathbb F_2=\{0,1\}\) for the field with two elements, and \(w(x)=|\{j:x_j=1\}|\) for Hamming weight. For positive integers \(n\) and integers \(1\le a\le b\), define
\[A(a,b,n)=\{x\in\mathbb F_2^n:a\le w(x)\le b\},\qquad
R_A(a,b,n)=\frac1n\max\{\dim V:V\le\mathbb F_2^n,\ V\cap A(a,b,n)=\varnothing\}.\]
This definition also applies when \(b>n\). Equivalently, \(1-R_A(a,b,n)\) is the least normalized rank of a binary matrix with kernel avoiding the annulus. Thus \(R_A\) is the rate of the avoiding subspace, rather than the normalized rank of its parity-check matrix.
For fixed \(0\le\alpha<\beta\le1\), put \(a_n=\max\{1,\lfloor\alpha n\rfloor\}\), \(b_n=\lfloor\beta n\rfloor\), and \(N_n=n-a_n+1\); take sufficiently large \(n\) that \(a_n\le b_n\).
\paragraph{Statement recorded in the supplied source.}
Does the asymptotic annulus identity
\[R_A(a_n,b_n,n)=\frac{a_n-1}{n}+\frac{N_n}{n}R_A(1,b_n,N_n)+o(1)\quad(n\to\infty)\]
hold for every fixed \(0\le\alpha<\beta\le1\)? Equivalently, in the usual notation suppressing integer rounding,
\[R_A(\alpha n,\beta n,n)=\alpha+(1-\alpha)R_A(1,\beta n,(1-\alpha)n)+o(1).\]
Here \(o(1)\) denotes a quantity tending to zero for the fixed pair \((\alpha,\beta)\). \sref{2800702}{2}
\subsection{Short English statement}
Does the largest rate of a binary linear subspace avoiding a Hamming annulus asymptotically equal the rate obtained by using free coordinates below its inner radius and a code avoiding a punctured ball on the remaining coordinates?
\subsection{Sources}
\begin{description}
\item[\hypertarget{source:2800702:1}{S1}] ULAM.AI and the UnsolvedMath Contributors, UnsolvedMath: A Curated Collection of Open Mathematics Problems, record 2800702 (AMR-027-0702). CC BY 4.0. \url{https://huggingface.co/datasets/ulamai/UnsolvedMath}.
\item[\hypertarget{source:2800702:2}{S2}] Emmanuel Abbe, Noga Alon, Afonso S. Bandeira and Colin Sandon, Linear Boolean Classification, Coding and “the Critical Problem” (2015), arXiv:1401.6528, Definition 5, equations (21)–(23); Proposition 1, equation (26); Conjecture 3, equation (34). \url{https://arxiv.org/pdf/1401.6528}.
\end{description}
\label{end:2800702}
\end{document}
